\section{Overlap Protection Handling}
\label{sec:interlocking:overlap-protection}

Overlaps are reserved track segments positioned beyond a signal, serving as a critical safety buffer should a train overshoot its intended stopping point. The system enforces overlap protection through both static and dynamic modes. In static configurations, each route is assigned a predefined overlap that must be proven unoccupied and securely locked before the signal can clear. In dynamic mode, the interlocking engine selects from a list of candidate overlaps and reserves the first available option, ensuring flexibility in complex station layouts. Once granted, an overlap remains locked until the train has fully cleared it, as confirmed by occupancy sensors, after which it is released back into operation. If no valid overlap can be secured, the requested route is denied and the signal remains at red. By enforcing this mechanism, the system provides a crucial safeguard against accidents arising from emergency braking failures or wheel slippage, ensuring that even unexpected overruns remain contained within a controlled and protected zone.
